\setcounter{dang}{0}
\section{Phương trình quy về phương trình bậc hai}
\subsection{Tóm tắt lý thuyết}
\subsubsection{Phương trình dạng $\sqrt{a x^{2}+b x+c}=\sqrt{d x^{2}+e x+f}$}
Để giải phương trình $\sqrt{a x^{2}+b x+c}=\sqrt{d x^{2}+e x+f}$, bình phương hai vế sau đó thu gọn ta được phương trình
$$
(a-d) x^{2}+(b-e) x+(c-f)=0. \quad (1)
$$
Giải phương trình (1) được các nghiệm, sau đó thay vào phương trình ban đầu để thử lại xem nghiệm nào thoả mãn và kết luận.

Chú ý rằng nếu $x_{0}$ là một nghiệm của phương trình (1) thì khi thử lại ta chỉ cần kiểm tra xem, nếu $a x_{0}^{2}+b x_{0}+c \geq 0$ thì $x_{0}$ sẽ là nghiệm của phương trình đã cho.
\subsubsection{Phương trình dạng $\sqrt{a x^{2}+b x+c}=d x+e$}
Để giải phương trình $\sqrt{a x^{2}+b x+c}=d x+e$, bình phương hai vế sau đó thu gọn ta được phương trình
$$
\left(a-d^{2}\right) x^{2}+(b-2 d e) x+\left(c-e^{2}\right)=0. \quad (2)
$$
Giải phương trình (2) được các nghiệm, sau đó thay vào phương trình ban đầu để thử lại xem nghiệm nào thoả mãn và kết luận.

Chú ý rằng nếu $x_{0}$ là một nghiệm của phương trình (2) thì khi thử lại ta chỉ cần kiểm tra xem, nếu $d x_{0}+e \geq 0$ thì $x_{0}$ sẽ là nghiệm của phương trình đã cho.


\subsection{Các dạng toán}
\begin{dang}{Phương trình $\sqrt{f(x)}=\sqrt{g(x)}$}
	\begin{itemize}
		\item Bình phương 2 vế và giải.
		\item Kiểm tra nghiệm và kết luận.
	\end{itemize}
\end{dang}
\viduminhhoa
\begin{vd}%[0D3B2-4]
	Giải phương trình $\sqrt{x^{2}-2x-4}=\sqrt{2-x}$.
	\loigiai{
		\allowdisplaybreaks
		\begin{eqnarray*}
			\sqrt{x^2-2x-4}=\sqrt{2-x} &\Rightarrow &
			x^2-2x-4=2-x\\
			&\Rightarrow&
			x^2-x-6=0\\
			&	\Rightarrow&\hoac{&x=-2\\&x=3.}
		\end{eqnarray*}
		Chỉ có $x=-2$ thỏa mãn phương trình ban đầu. \\ Vậy phương trình có nghiệm $x=-2$.}
\end{vd}
\begin{vd}%[0D3B2-4]
	Giải phương trình $\sqrt{2x^2-5x+2}-\sqrt{6-3x}=0$.
	\loigiai{
		\allowdisplaybreaks
		\begin{eqnarray*}
			\sqrt{2x^2-5x+2}-\sqrt{6-3x}=0 
			&\Rightarrow & \sqrt{2x^2-5x+2}=\sqrt{6-3x}\\
			&\Rightarrow &  2x^2-5x+2 = 6-3x\\
			& \Rightarrow &  x^2-x-2 = 0\\
			&\Rightarrow &  \hoac{& x = -1 \\ & x = 2.}
		\end{eqnarray*}
		Các nghiệm thỏa mãn phương trình ban đầu. \\ Vậy nghiệm của phương trình là $x=-1$ và $x=2$.
	}
\end{vd}
\begin{vd}%[0D3B2-4]
	Giải phương trình $\sqrt{x^2+3x}=2\sqrt{3x-2}$.
	\loigiai{
		\allowdisplaybreaks
		\begin{eqnarray*}
			\sqrt{x^2+3x}=2\sqrt{3x-2}
			&\Rightarrow &  x^2 + 3x = 4(3x - 2) \\
			&\Rightarrow & x^2 - 9x + 8 = 0\\
			&\Rightarrow & \hoac{& x = 1 \\ & x = 8.}
		\end{eqnarray*}
		Các nghiệm thỏa mãn phương trình ban đầu. \\ Vậy nghiệm của phương trình là $x=1$ và $x=8$.
	}
\end{vd}
\begin{vd}%[0D3G2-4]
	Tìm điều kiện của tham số $m$ để phương trình sau có nghiệm:
	$$
	\sqrt{x^{2}+x+1}=\sqrt{2 x^{2}+m x+m+1}. \quad \hfill (1)
	$$
	\loigiai{
		Bình phương hai vế của phương trình (1) và thu gọn ta được
		$$
		x^{2}+(m-1) x+m=0. \quad (2)
		$$
		Nhận thấy rằng tam thức bậc hai $x^{2}+x+1$ có $a=1>0$ và $\Delta=-3<0$. \\
		Suy ra $x^{2}+x+1>0$ với mọi $x$. Như vậy nếu phương trình (2) có nghiệm $x_{0}$ thì khi thử lại ta thấy $x_{0}^{2}+x_{0}+1>0$, tức là $x_{0}$ thoả mãn phương trình (1).\\
		Vậy phương trình (1) có nghiệm khi và chỉ khi phương trình (2) có nghiệm.\\
		Điều này tương đương với $\Delta=(m-1)^{2}-4 m \geq 0$ hay $m^{2}-6 m+1 \geq 0$.\\
		Từ đó ta được $m \leq 3-2 \sqrt{2}$ hoặc $m \geq 3+2 \sqrt{2}$.
	}
\end{vd}
\baitaptl
\begin{bt}%[0D3K2-4]
	Giải phương trình  $\sqrt{3x^{2}-4x-4}=\sqrt{2x+5}$.
	\loigiai{\allowdisplaybreaks
		\begin{eqnarray*}
			\sqrt{3x^2-4x-4}=\sqrt{2x+5} &\Rightarrow&
			3x^2-4x-4=2x+5\\
			&\Rightarrow& \hoac{&x=-1\\
				&x=3.}
		\end{eqnarray*}
		Các nghiệm thỏa mãn phương trình ban đầu. \\ Vậy phương trình có nghiệm $x=-1; x=3$.}
\end{bt}
\begin{bt}%[0D3B2-4] 
	Giải phương trình $\sqrt{2x-1}-\dfrac{1}{2}\sqrt{x+3}=0$.
	\loigiai{
		\allowdisplaybreaks
		\begin{eqnarray*}
			\sqrt{2x-1}-\dfrac{1}{2}\sqrt{x+3}=0 
			&	\Rightarrow &\sqrt{2x-1}=\dfrac{1}{2}\sqrt{x+3}\\
			&	\Rightarrow & 2x-1 = \dfrac{1}{4}(x+3)\\
			& \Rightarrow & 4(2x-1) = x+3 \\
			&\Rightarrow & 7x -7 = 0\\
			&	\Rightarrow & x = 1.
		\end{eqnarray*}
		Nghiệm thỏa mãn phương trình ban đầu. \\ Vậy nghiệm của phương trình là $x=1$.}
\end{bt}
\begin{bt}%[0D3B2-4] 
	Giải phương trình  $\sqrt{x^2+3x}=\sqrt{2(x+1)}$.
	\loigiai{	
		\allowdisplaybreaks
		\begin{eqnarray*}
			\sqrt{x^2+3x}=\sqrt{2(x+1)} 
			&\Rightarrow  & x^2+3x = 2(x+1) \\
			&\Rightarrow& x^2+x-2 = 0 \\
			&	\Rightarrow&  \hoac{& x = 1\\ & x = -2.}
		\end{eqnarray*}
		Chỉ có nghiệm $x = 1$ thỏa mãn phương trình ban đầu. \\ Vậy nghiệm của phương trình là $x=1$.}
\end{bt}
\begin{bt}%[0D3B2-4] 
	Giải phương trình $\sqrt{2x^2-3x+4}=\dfrac{3}{2}\sqrt{x+5}$.
	\loigiai{
		\allowdisplaybreaks
		\begin{eqnarray*}
			\sqrt{2x^2-3x+4}=\dfrac{3}{2}\sqrt{x+5} 
			&	\Rightarrow &2x^2-3x+4 = \dfrac{9}{4}(x+5) \\
			&\Rightarrow &  4(2x^2-3x+4) = 9(x+5)\\
			& \Rightarrow &  8x^2-21x-29 = 0 \\
			&\Rightarrow &\hoac{& x = -1\\ & x = \dfrac{29}{8}.}
		\end{eqnarray*}	
		Các nghiệm thỏa mãn phương trình ban đầu. \\ Vậy nghiệm của phương trình là $x=-1$ và $x = \dfrac{29}{8}$.}
\end{bt}
\begin{bt}%[0D3B2-4]
	Giải phương trình sau:
	$\sqrt{2x^2+x}-\sqrt{3}=0$.
	\loigiai{
		\allowdisplaybreaks
		\begin{eqnarray*}
			\sqrt{2x^2+x}-\sqrt{3}=0 
			&\Rightarrow &\sqrt{2x^2+x}=\sqrt{3} \\
			&	\Rightarrow & 2x^2 + x = 3\\
			&	\Rightarrow & 2x^2 + x -3 = 0 \\
			&	\Rightarrow & \hoac{& x = 1 \\ & x = -\dfrac{3}{2}.}
		\end{eqnarray*}
		Các nghiệm thỏa mãn phương trình ban đầu. \\ Vậy nghiệm của phương trình là $x=1$ và $x = -\dfrac{3}{2}$.
	}
\end{bt}
\begin{bt}%[0D3B2-4]
	Giải phương trình sau:
	$2\sqrt{(x+1)^2+7}=\sqrt{7(x+5)}$.
	\loigiai{
		\allowdisplaybreaks
		\begin{eqnarray*}
			2\sqrt{(x+1)^2+7}=\sqrt{7(x+5)} 
			&	\Rightarrow&  4[(x+1)^2+7] = 7(x+5) \\
			&	\Rightarrow& 4(x^2+2x+8) = 7x+35 \\
			& \Rightarrow & 4x^2 + x - 3 = 0 \\
			&	\Rightarrow & \hoac{& x = -1  \\ & x = \dfrac{3}{4}.}
		\end{eqnarray*}
		Các nghiệm thỏa mãn phương trình ban đầu. \\ Vậy nghiệm của phương trình là $x=-1$ và $x = \dfrac{3}{4}$.
	}
\end{bt}
\begin{bt}%[0D3B2-4]
	Giải  phương trình sau:
	$\sqrt{(x+1)(x-3)}-\sqrt{x}=0$.
	\loigiai{
		\allowdisplaybreaks
		\begin{eqnarray*}
			\sqrt{(x+1)(x-3)}-\sqrt{x}=0 
			&\Rightarrow& \sqrt{(x+1)(x-3)}=\sqrt{x} \\
			&\Rightarrow &  (x+1)(x-3) = x \\
			& \Rightarrow &  x^2 -3x -3 = 0\\ 
			&\Rightarrow& \hoac{& x = \dfrac{3 + \sqrt{21}}{2} \\ & x = \dfrac{3 - \sqrt{21}}{2}.}
		\end{eqnarray*}	
		Chỉ có nghiệm $x = \dfrac{3 + \sqrt{21}}{2}$ thỏa mãn phương trình ban đầu. \\ Vậy nghiệm của phương trình là $x = \dfrac{3 + \sqrt{21}}{2}$.	
	}
\end{bt}

\begin{bt}%[0D3B2-4]
	Giải phương trình $\sqrt{x - 3} = 3\sqrt{x^2 -9}$.
	\loigiai{
		\begin{eqnarray*} \sqrt{x - 3} = 3\sqrt{x^2 -9}
			&\Rightarrow& x-3 = 9\left(x^2-9\right)\\
			& \Rightarrow& 9x^2 - x -78= 0\\
			&\Rightarrow& \hoac{&x =  3\\&x = -\dfrac{26}{9}.}
		\end{eqnarray*}	
		Chỉ có nghiệm $x=3$ thỏa mãn phương trình ban đầu. \\ Vậy phương trình có tập nghiệm $S = \left\{3\right\}$.
	}
\end{bt}

\begin{bt}%[0D3K2-4]
	Giải phương trình $\sqrt{x^2 - 1+\sqrt{2x+3}} = \sqrt{x^2 + x-1}$.
	\loigiai{
		\begin{eqnarray*} \sqrt{x^2 - 1+\sqrt{2x+3}} = \sqrt{x^2 + x-1}
			&\Rightarrow& x^2 - 1+\sqrt{2x+3} = x^2 + x-1\\
			& \Rightarrow& \sqrt{2x+3}= x\\
			& \Rightarrow&  2x+3 = x^2\\
			& \Rightarrow & x^2-2x-3=0\\
			& \Rightarrow & \hoac{&x =  3\\& x= -1.}
		\end{eqnarray*}
		Chỉ có nghiệm $x=3$ thỏa mãn phương trình ban đầu. \\ Vậy phương trình có tập nghiệm $S = \left\{3\right\} $.
	}
\end{bt}
\begin{bt}%[0D3G2-3]
	Tìm điều kiện của tham số $m$ để phương trình sau có nghiệm:
	$$
	\sqrt{2 x^{2}+x+1}=\sqrt{x^{2}+m x+m-1}. \quad \hfill (1)
	$$
	\loigiai{
		Bình phương hai vế phương trình đã cho và thu gọn ta được 
		$$x^{2}+(1-m) x+2-m=0. \quad \hfill(2)$$
		Do $2 x^{2}+x+1>0$ với mọi $x \in \mathbb{R}$ (vì $a=2>0$ và $\Delta=-7<0$ ) nên nếu $x_{0}$ là một nghiệm của phương trình (2) thì khi thử lại ta được $2 x_{0}^{2}+x_{0}+1>0$, tức là $x_{0}$ thoả mãn phương trình đã cho.\\
		Vậy phương trình đã cho có nghiệm khi và chỉ khi phương trình (2) có nghiệm.\\
		Điều này tương đương với
		\begin{eqnarray*}
			&&\Delta=(1-m)^{2}-4(2-m) \geq 0 \\
			& \Leftrightarrow & m^{2}+2 m-7 \geq 0 \\
			& \Leftrightarrow & m \geq-1+2 \sqrt{2} \text { hoặc } m \leq-1-2 \sqrt{2}.
		\end{eqnarray*}
	}
\end{bt}
\subsubsection{Câu hỏi trắc nghiệm}
\setcounter{ex}{0}
\Opensolutionfile{ans}[ans/ans-18-dang-1]
\begin{ex}%[0D3B2-3]
	Phương trình sau có bao nhiêu nghiệm $\sqrt{x}=\sqrt{-x}$ 
	\choice
	{$0$}
	{\True $1$}
	{$2$}
	{Vô số}
	\loigiai{
		Điều kiện xác định: $\heva{x &\geq 0 \\ -x &\geq 0}\Rightarrow x=0$. \\
		Ta thấy $x=0$ thoả mãn phương trình. Nghiệm thỏa mãn phương trình ban đầu. \\ Vậy phương trình đã cho có 1 nghiệm.}
\end{ex}

\begin{ex}%[0D3B2-1]
	Điều kiện xác định của phương trình $\sqrt{x-1}=\sqrt{x+2}$ là
	\choice
	{$x>1$}
	{$x \geq -2$}
	{$x>-2$}
	{\True $x \geq 1$}
	\loigiai{Do $x+2>x-1$ nên\\
		Phương trình xác định khi và chỉ khi $x-1 \geq 0 \Rightarrow x \geq 1$.}
\end{ex}
\begin{ex}%[0D3B2-4]
	Phương trình sau có bao nhiêu nghiệm: $\sqrt{x-2}=\sqrt{2-x}$
	\choice
	{$0$}
	{\True $1$}
	{$2$}
	{Vô số}
	\loigiai{
		Điều kiện $\heva{&x-2 \geq 0 \\& 2-x \geq 0}\Rightarrow \heva{&x \geq 2 \\& x \leq 2}\Rightarrow x=2$.\\
		Thay $x=2$ vào phương trình ta được $0=0$ hay $x=2$ là nghiệm của phương trình.}
\end{ex}
\begin{ex}%[0D3B2-4]
	Nghiệm của phương trình $\sqrt{x-1}=\left(\sqrt{3-x}\right)^2$ là
	\choice
	{$x=2;x=5$}
	{\True $x=2$}
	{$x=1;x=3$}
	{$x=-1;x=-3$}
	\loigiai{
		Ta có
		\allowdisplaybreaks
		\begin{eqnarray*}
			\sqrt{x-1}=\left(\sqrt{3-x}\right)^2	&\Rightarrow & \sqrt{x-1}=3-x\\
			& \Rightarrow  & x-1=(3-x)^2 \\
			& \Rightarrow & x^2-7x+10=0\\
			&\Rightarrow & \hoac{&x=2\\&x=5}.
		\end{eqnarray*}
		Chỉ có nghiệm $x=2$ thỏa mãn phương trình ban đầu. \\
		Vậy $S=\{2\}$.
	}
\end{ex}
\begin{ex}%[0D3B2-4]
	Tổng các nghiệm của phương trình $\sqrt{3x^2 - 8x + 5} - \sqrt{11 - x}=0 $ là
	\choice
	{\True $\dfrac{7}{3}$}
	{$\dfrac{11}{3}$}
	{$-\dfrac{11}{3}$}
	{$\dfrac{1}{3}$}
	\loigiai{
		\begin{eqnarray*} \text{Phương trình} 
			&\Rightarrow \sqrt{3x^2 - 8x + 5} = \sqrt{11 - x}\\
			& \Rightarrow 3x^2 - 8x + 5 = 11 - x\\
			& \Rightarrow 3x^2 - 7x -6 = 0\\
			&\Rightarrow \hoac{&x =  3\\&x = -\dfrac{2}{3}.}\\
		\end{eqnarray*}	
		Các nghiệm thỏa mãn phương trình ban đầu. \\
		Nên phương trình có tập nghiệm $S = \left\lbrace3; -\dfrac{2}{3} \right\rbrace $.\\
		Vậy tổng các nghiệm là $\dfrac{7}{3}$.
	}
\end{ex}
\begin{ex}%[0D3B2-4]
	Phương trình 
	$\sqrt{2x^2-3x+12}=2\sqrt{-x^2+x+3}$ có các nghiệm 
	\choice
	{\True không âm}
	{không dương}
	{dương}
	{âm}
	\loigiai{\begin{eqnarray*} \sqrt{2x^2-3x+12}=2\sqrt{-x^2+x+3}
			&\Rightarrow& 2x^2-3x+12 = 4\left(-x^2+x+3\right)\\
			& \Rightarrow & 6x^2 -7x= 0\\
			& \Rightarrow & \hoac{& x=0  \\& x=\dfrac{7}{6}.}
		\end{eqnarray*}	
		Các nghiệm thỏa mãn phương trình ban đầu. \\ Vậy phương trình có  các nghiệm không âm và tập nghiệm $S = \left\{ 0; \dfrac{7}{6} \right\}$.  }
\end{ex}
\begin{ex}%[0D3B2-4]
	Số nghiệm của phương trình 	$\sqrt{x^2-3x-2}=\sqrt{x-3}$ là
	\choice
	{\True 1}
	{0}
	{2}
	{3}
	\loigiai{
		\begin{eqnarray*} \sqrt{x^2-3x-2}=\sqrt{x-3}
			&\Rightarrow& x^2-3x-2 = x-3\\
			& \Rightarrow & x^2 -4x+1= 0\\
			& \Rightarrow & \hoac{& x=2+\sqrt{3} \\& x=2-\sqrt{3}.}
		\end{eqnarray*}
		Chỉ có nghiệm $x=2+\sqrt{3}$ thỏa mãn phương trình ban đầu. \\ Vậy phương trình có 1 nghiệm và $S=\{2+\sqrt{3}\} $
	}
\end{ex}	
\begin{ex}%[0D3B2-4]
	Số nghiệm của phương trình	$\sqrt{x^2-3x+18}=\sqrt{14x+2}$ là
	\choice
	{\True 2}
	{0}
	{1}
	{3}
	\loigiai{Điều kiện $\heva{&x^2-3x+18 \ge 0\\&14x+2 \ge0.}$
		\begin{eqnarray*} \sqrt{x^2-3x+18}=\sqrt{14x+2} 
			&\Rightarrow & x^2-3x+18 = 14x+2\\
			& \Rightarrow & x^2 -17x+16= 0\\
			& \Rightarrow & \hoac{& x=1  \\& x=16.}
		\end{eqnarray*}
		Các nghiệm thỏa mãn phương trình ban đầu. \\ Vậy phương trình có 2 nghiệm và tập nghiệm $S = \left\{1;16\right\} $.  }
\end{ex}
\begin{ex}%[0D3B2-4]
	Số nghiệm của phương trình	$\sqrt{x^2-5x+2}=\sqrt{-x-1}$ là 
	\choice
	{\True 0}
	{1}
	{2}
	{4}
	\loigiai{\begin{eqnarray*} \sqrt{x^2-5x+2}=\sqrt{-x-1}
			&\Rightarrow & x^2-5x+2 = -x-1\\
			& \Rightarrow & x^2 -4x+3= 0\\
			& \Rightarrow & \hoac{& x=1  \\& x=3.}
		\end{eqnarray*}	
		Các nghiệm không thỏa mãn phương trình ban đầu. \\ Vậy phương trình có tập nghiệm $S =\varnothing $.}
\end{ex}
\begin{ex}%[0D3B2-4]
	Tổng bình phương các nghiệm của	 phương trình $3\sqrt{x-1}=\sqrt{x^2+8x-11}$ là
	\choice
	{\True 4}
	{8}
	{5}
	{7}
	\loigiai{\begin{eqnarray*} 3\sqrt{x-1}=\sqrt{x^2+8x-11}
			&\Rightarrow& 9\left(x-1\right) = x^2+8x-11\\
			& \Rightarrow& x^2 -x-2= 0\\
			& \Rightarrow & \hoac{& x=2  \\& x=-1.}
		\end{eqnarray*}
		Chỉ có nghiệm $x=2$ thỏa mãn phương trình ban đầu.\\
		Nên phương trình có tập nghiệm $S = \left\{ 2 \right\}$. \\
		Vậy tổng bình phương các nghiệm là $4$.}
\end{ex}
\Closesolutionfile{ans}

\begin{indapan}{10}
	{ans/ans-18-dang-1}
\end{indapan}
\begin{dang}{Phương trình $\sqrt{f(x)}=g(x)$}
	\begin{itemize}
		\item Giải bất phương trình $g(x)\geq 0\; (*)$ để tìm tập nghiệm của bất phương trình này.
		\item Bình phương hai vế của phương trình ta được $f(x)=\left[g(x)\right]^2$ rồi tìm tập nghiệm của phương trình đó.
		\item Trong những nghiệm của phương trình trên, ta giữ lại những nghiệm thuộc tập nghiệm của bất phương trình $(*)$. Tập nghiệm giữ lại đó chính là tập nghiệm cần tìm.
	\end{itemize}
\end{dang}
\viduminhhoa
\begin{vd}%[0D3B2-4]
	Giải phương trình $\sqrt{x^2+3x-1}=3$.
	\loigiai{
		\allowdisplaybreaks
		\begin{eqnarray*}
			&&\sqrt{x^2+3x-1}=3\\
			&\Rightarrow&x^2+3x-1=3^2\\
			&\Rightarrow&x^2+3x-10=0\\
			&\Rightarrow&\hoac{&x=2\\&x=-5.}
		\end{eqnarray*}
		Trong hai giá trị trên đều thỏa mãn.\\
		Vậy phương trình đã cho có tập nghiệm $S=\{2;-5\}$.
	}
\end{vd}

\begin{vd}%[0D3B2-4]
	Giải phương trình $\sqrt{x^2-8x+13}=2x-3$.
	\loigiai{
		Trước hết ta giải bất phương trình $2x-3\geq 0\Leftrightarrow x\geq \dfrac{3}{2}.\; (*)$
		\allowdisplaybreaks
		\begin{eqnarray*}
			&&\sqrt{x^2-8x+13}=2x-3\\
			&\Rightarrow&x^2-8x+13=(2x-3)^2\\
			&\Rightarrow&x^2-8x+13=4x^2-12x+9\\
			&\Rightarrow&3x^2-4x-4=0\\
			&\Rightarrow&\hoac{&x=2\\&x=-\dfrac{2}{3}.}
		\end{eqnarray*}
		Trong hai giá trị trên, chỉ có $x=2$ thỏa mãn $(*)$.\\
		Vậy phương trình đã cho có tập nghiệm $S=\{2\}$.
	}
\end{vd}
\begin{vd}%[0D3B2-4]
	Giải phương trình $\sqrt{-x^2-2x+1}=-x-3$.
	\loigiai{
		Trước hết ta giải bất phương trình $-x-3\geq 0\Leftrightarrow x\leq -3.\; (*)$
		\allowdisplaybreaks
		\begin{eqnarray*}
			&&\sqrt{-x^2-2x+1}=-x-3\\
			&\Rightarrow&-x^2-2x+1=(-x-3)^2\\
			&\Rightarrow&-x^2-2x+1=x^2+6x+9\\
			&\Rightarrow&2x^2+8x+8=0\\
			&\Rightarrow&x=-2.
		\end{eqnarray*}
		Thấy $x=-2$ không thỏa mãn $(*)$.\\
		Vậy phương trình đã cho vô nghiệm.
	}
\end{vd}

\begin{vd}%[0D3G2-4]
	Tìm điều kiện của tham số $m$ để phương trình sau có đúng một nghiệm:
	$$
	\sqrt{2x+m-4}=x-1. \quad \hfill (1)
	$$
	\loigiai{
		Trước hết ta giải bất phương trình $x-1\geq 0\Leftrightarrow x\geq 1.\qquad (*)$\\
		Bình phương hai vế của phương trình $(1)$ ta được $$2x+m-4=(x-1)^2\Leftrightarrow m=x^2-4x+5.\qquad (2)$$
		Phương trình $(1)$ có đúng một nghiệm khi và chỉ khi phương trình $(2)$ có đúng một nghiệm thỏa mãn $(*)$.\\
		Xét hàm số $f(x)=x^2-4x+5$ với $x\geq 1$.\\
		Bảng biến thiên
		\begin{center}
			\begin{tikzpicture}
				\tkzTabInit[nocadre=false,lgt=1.2,espcl=2.5,deltacl=0.6]
				{$x$ /0.6, $f(x)$ /2}
				{$1$,$2$,$+\infty$}
				\tkzTabVar{+/$2$,-/$1$,+/$+\infty$}
			\end{tikzpicture}
		\end{center}
		Từ bảng biến thiên suy ra, phương trình $(2)$ có đúng một nghiệm thỏa mãn $(*)$ khi  và chỉ khi $$\hoac{&m=1\\&m>2.}$$
		Vậy với $\hoac{&m=1\\&m>2.}$ thì phương trình $(1)$ có đúng một nghiệm.
	}
\end{vd}
\baitaptl
\setcounter{bt}{0}
\begin{bt}%[0D3B2-4]
	Giải phương trình $\sqrt{2-x^2}+x=0$.
	\loigiai{
		Ta có $\sqrt{2-x^2}+x=0\Leftrightarrow \sqrt{2-x^2}=-x$.\\
		Trước hết ta giải bất phương trình $-x\geq 0\Leftrightarrow x\leq 0.\; (*)$
		\allowdisplaybreaks
		\begin{eqnarray*}
			&&\sqrt{2-x^2}=-x\\
			&\Rightarrow&2-x^2=(-x)^2\\
			&\Rightarrow&x^2=1\\
			&\Rightarrow&\hoac{&x=1\\&x=-1.}
		\end{eqnarray*}
		Trong hai giá trị trên, chỉ có $x=-1$ thỏa mãn $(*)$.\\
		Vậy phương trình đã cho có tập nghiệm $S=\{-1\}$.
	}
\end{bt}
\begin{bt}%[0D3B2-4]
	Giải phương trình $\sqrt{3x-8}=x-2$.
	\loigiai{
		Trước hết ta giải bất phương trình $x-2\geq 0\Leftrightarrow x\geq 2.\; (*)$
		\allowdisplaybreaks
		\begin{eqnarray*}
			&&\sqrt{3x-8}=x-2\\
			&\Rightarrow&3x-8=(x-2)^2\\
			&\Rightarrow&x^2-7x+12=0\\
			&\Rightarrow&\hoac{&x=3\\&x=4.}
		\end{eqnarray*}
		Cả hai giá trị trên đều thỏa mãn $(*)$.\\
		Vậy phương trình đã cho có tập nghiệm $S=\{3;4\}$.
	}
\end{bt}
\begin{bt}%[0D3B2-4]
	Giải phương trình $\sqrt{5x-x^2}+x=3$.
	\loigiai{
		Ta có $\sqrt{5x-x^2}+x=3\Leftrightarrow \sqrt{5x-x^2}=3-x$.\\
		Trước hết ta giải bất phương trình $3-x\geq 0\Leftrightarrow x\leq 3.\; (*)$
		\allowdisplaybreaks
		\begin{eqnarray*}
			&&\sqrt{5x-x^2}=3-x\\
			&\Rightarrow&5x-x^2=(3-x)^2\\
			&\Rightarrow&2x^2-11x+9=0\\
			&\Rightarrow&\hoac{&x=1\\&x=\dfrac{9}{2}.}
		\end{eqnarray*}
		Trong hai giá trị trên, chỉ có $x=1$ thỏa mãn $(*)$.\\
		Vậy phương trình đã cho có tập nghiệm $S=\{1\}$.
	}
\end{bt}
\begin{bt}%[0D3B2-4]
	Giải phương trình $\sqrt{8x^2-10x+1}=2x+1$.
	\loigiai{
		Trước hết ta giải bất phương trình $2x+1\geq 0\Leftrightarrow x\geq -\dfrac{1}{2}.\; (*)$
		\allowdisplaybreaks
		\begin{eqnarray*}
			&&\sqrt{8x^2-10x+1}=2x+1\\
			&\Rightarrow&8x^2-10x+1=(2x+1)^2\\
			&\Rightarrow&4x^2-14x=0\\
			&\Rightarrow&\hoac{&x=0\\&x=\dfrac{7}{2}.}
		\end{eqnarray*}
		Cả hai giá trị trên đều thỏa mãn $(*)$.\\
		Vậy phương trình đã cho có tập nghiệm $S=\left\{0;\dfrac{7}{2}\right\}$.
	}
\end{bt}
\begin{bt}%[0D3B2-4]
	Giải phương trình $\sqrt{x^2-3x+2}=-1$.
	\loigiai{
		Ta thấy vế phải của phương trình luôn âm. Khi $\sqrt{x^2-3x+2}$ có nghĩa thì vế trái phương trình luôn không âm.\\
		Vậy phương trình đã cho vô nghiệm.
	}
\end{bt}
\begin{bt}%[0D3B2-4]
	Giải phương trình $\sqrt{3x^2+63x+229}=2x+15$.
	\loigiai{
		Trước hết ta giải bất phương trình $2x+15\geq 0\Leftrightarrow x\geq -\dfrac{15}{2}.\; (*)$
		\allowdisplaybreaks
		\begin{eqnarray*}
			&&\sqrt{3x^2+63x+229}=2x+15\\
			&\Rightarrow&3x^2+63x+229=(2x+15)^2\\
			&\Rightarrow&x^2-3x-4=0\\
			&\Rightarrow&\hoac{&x=-1\\&x=4.}
		\end{eqnarray*}
		Cả hai giá trị trên đều thỏa mãn $(*)$.\\
		Vậy phương trình đã cho có tập nghiệm $S=\left\{-1;4\right\}$.
	}
\end{bt}
\begin{bt}%[0D3B2-4]
	Giải phương trình $\sqrt{x^2+7x+11}=2$.
	\loigiai{
		\allowdisplaybreaks
		\begin{eqnarray*}
			&&\sqrt{x^2+7x+11}=2\\
			&\Rightarrow&x^2+7x+11=2^2\\
			&\Rightarrow&x^2+7x+7=0\\
			&\Rightarrow&\hoac{&x=\dfrac{-7+\sqrt{21}}{2}\\&x=\dfrac{-7-\sqrt{21}}{2}.}
		\end{eqnarray*}
		Trong hai giá trị trên đều thỏa mãn.\\
		Vậy phương trình đã cho có tập nghiệm $S=\left\{\dfrac{-7+\sqrt{21}}{2};\dfrac{-7-\sqrt{21}}{2}\right\}$.
	}
\end{bt}
\begin{bt}%[0D3B2-4]
	Giải phương trình $1+\sqrt{11-x}=x$.
	\loigiai{
		Ta có $1+\sqrt{11-x}=x\Leftrightarrow \sqrt{11-x}=x-1$.\\
		Trước hết ta giải bất phương trình $x-1\geq 0\Leftrightarrow x\geq 1.\; (*)$
		\allowdisplaybreaks
		\begin{eqnarray*}
			&&\sqrt{11-x}=x-1\\
			&\Rightarrow&11-x=(x-1)^2\\
			&\Rightarrow&x^2-x-10=0\\
			&\Rightarrow&\hoac{&x=\dfrac{1+\sqrt{41}}{2}\\&x=\dfrac{1-\sqrt{41}}{2}.}
		\end{eqnarray*}
		Trong hai giá trị trên, chỉ có $x=\dfrac{1+\sqrt{41}}{2}$ thỏa mãn $(*)$.\\
		Vậy phương trình đã cho có tập nghiệm $S=\left\{\dfrac{1+\sqrt{41}}{2} \right\}$.
	}
\end{bt}
\begin{bt}%[0D3B2-4]
	Giải phương trình $\sqrt{2x^2+x}=-x-2$.
	\loigiai{
		Trước hết ta giải bất phương trình $-x-2\geq 0\Leftrightarrow x\leq -2.\; (*)$
		\allowdisplaybreaks
		\begin{eqnarray*}
			&&\sqrt{2x^2+x}=-x-2\\
			&\Rightarrow&2x^2+x=(-x-2)^2\\
			&\Rightarrow&x^2-3x-4=0\\
			&\Rightarrow&\hoac{&x=-1\\&x=4.}
		\end{eqnarray*}
		Cả hai giá trị trên đều không thỏa mãn $(*)$.\\
		Vậy phương trình đã cho vô nghiệm.
	}
\end{bt}

\begin{bt}%[0D3G2-4]
	Tìm điều kiện của tham số $m$ để phương trình sau có hai nghiệm phân biệt:
	$$
	\sqrt{-x^2+6x+m}=x-1. \quad \hfill (1)
	$$
	\loigiai{
		Trước hết ta giải bất phương trình $x-1\geq 0\Leftrightarrow x\geq 1.\qquad (*)$\\
		Bình phương hai vế của phương trình $(1)$ ta được $$-x^2+6x+m=(x-1)^2\Leftrightarrow m=2x^2-8x+1.\qquad (2)$$
		Phương trình $(1)$ có hai nghiệm phân biệt khi và chỉ khi phương trình $(2)$ có hai nghiệm phân biệt thỏa mãn $(*)$.\\
		Xét hàm số $f(x)=2x^2-8x+1$ với $x\geq 1$.\\
		Bảng biến thiên
		\begin{center}
			\begin{tikzpicture}
				\tkzTabInit[nocadre=false,lgt=1.2,espcl=2.5,deltacl=0.6]
				{$x$ /0.6, $f(x)$ /2}
				{$1$,$2$,$+\infty$}
				\tkzTabVar{+/$-5$,-/$-7$,+/$+\infty$}
			\end{tikzpicture}
		\end{center}
		Từ bảng biến thiên suy ra, phương trình $(2)$ có hai nghiệm phân biệt thỏa mãn $(*)$ khi và chỉ khi $-7<m<-5$.\\
		Vậy với $-7<m<-5$ thì phương trình $(1)$ có hai nghiệm phân biệt.
	}
\end{bt}

\subsubsection{Câu hỏi trắc nghiệm}
\setcounter{ex}{0}
\Opensolutionfile{ans}[ans/ans-18-dang-2]
\begin{ex}%[0D3B2-4]
	Số nghiệm của phương trình $\sqrt{4x^2-5x+1}=-3$ là
	\choice
	{\True $0$}
	{$1$}
	{$2$}
	{vô số}
	\loigiai{
		Ta thấy vế phải của phương trình luôn âm. Khi $\sqrt{4x^2-5x+1}$ có nghĩa thì vế trái phương trình luôn không âm.\\
		Vậy phương trình đã cho vô nghiệm.
	}
\end{ex}
\begin{ex}%[0D3B2-4]
	Tập nghiệm của phương trình $\sqrt{2x^2-3}=x-1$ là
	\choice
	{$\left\{-1-\sqrt{5};-1+\sqrt{5}\right\}$}
	{$\left\{-1-\sqrt{5}\right\}$}
	{\True $\left\{-1+\sqrt{5}\right\}$}
	{$\varnothing$}
	\loigiai{
		Trước hết ta giải bất phương trình $x-1\geq 0\Leftrightarrow x\geq 1.\; (*)$
		\allowdisplaybreaks
		\begin{eqnarray*}
			&&\sqrt{2x^2-3}=x-1\\
			&\Rightarrow&2x^2-3=(x-1)^2\\
			&\Rightarrow&x^2+2x-4=0\\
			&\Rightarrow&\hoac{&x=-1-\sqrt{5}\\&x=-1+\sqrt{5}.}
		\end{eqnarray*}
		Trong hai giá trị trên, chỉ có $x=-1+\sqrt{5}$ thỏa mãn $(*)$.\\
		Vậy phương trình đã cho có tập nghiệm $S=\left\{-1+\sqrt{5}\right\}$.
	}
\end{ex}

\begin{ex}%[0D3B2-4]
	Tập nghiệm của phương trình $\sqrt{2x^2+9x+10}=-x-4$ là
	\choice
	{$\left\{-3;2\right\}$}
	{$\left\{2\right\}$}
	{$\left\{-3\right\}$}
	{\True $\varnothing$}
	\loigiai{
		Trước hết ta giải bất phương trình $-x-4\geq 0\Leftrightarrow x\leq -4.\; (*)$
		\allowdisplaybreaks
		\begin{eqnarray*}
			&&\sqrt{2x^2+9x+10}=-x-4\\
			&\Rightarrow&2x^2+9x+10=(-x-4)^2\\
			&\Rightarrow&x^2+x-6=0\\
			&\Rightarrow&\hoac{&x=-3\\&x=2.}
		\end{eqnarray*}
		Trong hai giá trị trên đều không thỏa mãn $(*)$.\\
		Vậy phương trình đã cho có tập nghiệm $S=\varnothing$.
	}
\end{ex}
\begin{ex}%[0D3B2-4]
	Tổng các nghiệm của phương trình $\sqrt{3x^2+23x+29}=x+4$ là
	\choice
	{\True $-1$}
	{$-\dfrac{15}{2}$}
	{$-\dfrac{13}{2}$}
	{$0$}
	\loigiai{
		Trước hết ta giải bất phương trình $x+4\geq 0\Leftrightarrow x\geq -4.\; (*)$
		\allowdisplaybreaks
		\begin{eqnarray*}
			&&\sqrt{3x^2+23x+29}=x+4\\
			&\Rightarrow&3x^2+23x+29=(x+4)^2\\
			&\Rightarrow&2x^2+15x+13=0\\
			&\Rightarrow&\hoac{&x=-1\\&x=-\dfrac{13}{2}.}
		\end{eqnarray*}
		Trong hai giá trị trên, chỉ có $x=-1$ thỏa mãn $(*)$.\\
		Vậy phương trình đã cho có tập nghiệm $S=\left\{-1\right\}$, nên tổng của chúng bằng $-1$.
	}
\end{ex}

\begin{ex}%[0D3B2-4]
	Tổng các nghiệm của phương trình $\sqrt{3x^2+45x+67}=2x+9$ là
	\choice
	{$2$}
	{$7$}
	{\True $9$}
	{$0$}
	\loigiai{
		Trước hết ta giải bất phương trình $2x+9\geq 0\Leftrightarrow x\geq -\dfrac{9}{2}.\; (*)$
		\allowdisplaybreaks
		\begin{eqnarray*}
			&&\sqrt{3x^2+45x+67}=2x+9\\
			&\Rightarrow&3x^2+45x+67=(2x+9)^2\\
			&\Rightarrow&x^2-9x+14=0\\
			&\Rightarrow&\hoac{&x=2\\&x=7.}
		\end{eqnarray*}
		Cả hai giá trị trên đều thỏa mãn $(*)$.\\
		Vậy phương trình đã cho có tập nghiệm $S=\left\{2;7\right\}$, nên tổng của chúng bằng $14$.
	}
\end{ex}
\begin{ex}%[0D3B2-4]
	Tập nghiệm của phương trình $\sqrt{x^2-3x-2}=x+1$ là
	\choice
	{$\left\{-1-\sqrt{5};-1+\sqrt{5}\right\}$}
	{$\left\{-1-\sqrt{5}\right\}$}
	{\True $\left\{-\dfrac{3}{5}\right\}$}
	{$\varnothing$}
	\loigiai{
		Trước hết ta giải bất phương trình $x+1\geq 0\Leftrightarrow x\geq -1.\; (*)$
		\allowdisplaybreaks
		\begin{eqnarray*}
			&&\sqrt{x^2-3x-2}=x+1\\
			&\Rightarrow&x^2-3x-2=(x+1)^2\\
			&\Rightarrow&5x=-3\\
			&\Rightarrow&x=-\dfrac{3}{5}.
		\end{eqnarray*}
		Giá trị $x=-\dfrac{3}{5}$ thỏa mãn $(*)$.\\
		Vậy phương trình đã cho có tập nghiệm $S=\left\{-\dfrac{3}{5}\right\}$.
	}
\end{ex}
\begin{ex}%[0D3B2-4]
	Biết $\alpha$ là nghiệm của phương trình $\sqrt{2x^2-4x-1}=2-x$. Chọn khẳng định đúng.
	\choice
	{$2<\alpha<3$}
	{$-2<\alpha<2$}
	{$\alpha^2>5$}
	{\True $-3<\alpha<-2$}
	\loigiai{
		Trước hết ta giải bất phương trình $2-x\geq 0\Leftrightarrow x\leq 2.\; (*)$
		\allowdisplaybreaks
		\begin{eqnarray*}
			&&\sqrt{2x^2-4x-1}=2-x=2-x\\
			&\Rightarrow&2x^2-4x-1=(2-x)^2\\
			&\Rightarrow&x^2=5\\
			&\Rightarrow&\hoac{&x=-\sqrt{5}\\&x=\sqrt{5}.}
		\end{eqnarray*}
		Trong hai giá trị trên, chỉ có $x=-\sqrt{5}$ thỏa mãn $(*)$.\\
		Vậy phương trình đã cho có một nghiệm $x=-\sqrt{5}$.
	}
\end{ex}

\begin{ex}%[0D3B2-4]
	Số nghiệm của phương trình $\sqrt{2x^2-8x-5}=5-3x$ là
	\choice
	{$2$}
	{\True $0$}
	{$1$}
	{vô số}
	\loigiai{
		Trước hết ta giải bất phương trình $5-3x\geq 0\Leftrightarrow x\leq \dfrac{5}{3}.\; (*)$
		\allowdisplaybreaks
		\begin{eqnarray*}
			&&\sqrt{2x^2-8x-5}=5-3x\\
			&\Rightarrow&2x^2-8x-5=(5-3x)^2\\
			&\Rightarrow&7x^2-22x+30=0. \; \text{ phương trình này vô nghiệm.}
		\end{eqnarray*}
		Vậy phương trình đã cho vô nghiệm.
	}
\end{ex}

\begin{ex}%[0D3K2-4]
	Có bao nhiêu giá trị nguyên của tham số $m\in[-10;0]$ để phương trình sau có nghiệm? $$\sqrt{x^2-x+m}=x-2.\qquad (1)$$
	\choice
	{$10$}
	{\True $9$}
	{$11$}
	{$0$}
	\loigiai{
		Trước hết ta giải bất phương trình $x-2\geq 0\Leftrightarrow x\geq 2.\qquad (*)$\\
		Bình phương hai vế của phương trình $(1)$ ta được $$x^2-x+m=(x-2)^2\Leftrightarrow x=\dfrac{4-m}{3}.\qquad (2)$$
		Phương trình $(1)$ có nghiệm khi và chỉ khi phương trình $(2)$ có nghiệm thỏa mãn $$x\geq 2\Leftrightarrow \dfrac{4-m}{3}\geq 2\Leftrightarrow m\leq -2.$$
		Lại có $m$ nguyên thuộc đoạn $[-10;0]$ nên $m\in\{-10;-9;\ldots -2\}$.
	}
\end{ex}
\begin{ex}%[0D3K2-4]
	Tất cả các giá trị của tham số $m$ để phương trình $$\sqrt{-x^2-4x+m}=x-1.\qquad (1)$$ có nghiệm là
	\choice
	{$m\geq \dfrac{1}{2}$}
	{$m>5$}
	{\True $m\geq 5$}
	{$\dfrac{1}{2}\leq m\leq 5$}
	\loigiai{
		Trước hết ta giải bất phương trình $x-1\geq 0\Leftrightarrow x\geq 1.\qquad (*)$\\
		Bình phương hai vế của phương trình $(1)$ ta được $$-x^2-4x+m=(x-1)^2\Leftrightarrow m=2x^2+2x+1.\qquad (2)$$
		Phương trình $(1)$ có nghiệm khi và chỉ khi phương trình $(2)$ có nghiệm thỏa mãn $x\geq 1$.\\
		Xét hàm số $f(x)=2x^2+2x+1$ với $x\geq 1$.\\
		Bảng biến thiên
		\begin{center}
			\begin{tikzpicture}
				\tkzTabInit[nocadre=false,lgt=1.2,espcl=2.5,deltacl=0.6]
				{$x$ /0.6, $f(x)$ /2}
				{$1$,$+\infty$}
				%\tkzTabLine{ ,+, }
				\tkzTabVar{-/$5$,+/$+\infty$}
			\end{tikzpicture}
		\end{center}
		Từ bảng biến thiên suy ra phương trình $(2)$ có nghiệm thỏa mãn $x\geq 1$ khi và chỉ khi $m\geq 5$.\\
	}
\end{ex}
\Closesolutionfile{ans}

\begin{indapan}{10}
	{ans/ans-18-dang-2}
\end{indapan}
\begin{dang}{Toán thực tế}
	
\end{dang}

\viduminhhoa

\begin{vd}%[0D4B4]
	Bác Việt sống và làm việc tại trạm hải đăng cách bờ biển $4$ km. Hằng tuần bác chèo thuyền vào vị trí gần nhất trên bờ biển là bến Bính để nhận hàng hoá do cơ quan cung cấp. Tuần này, do trục trặc về vận chuyển nên toàn bộ số hàng vẫn đang nằm ở thôn Hoành, bên bờ biển cách bến Bính $9{,}25$ km và sẽ được anh Nam vận chuyển trên con đường dọc bờ biển tới bến Bính bằng xe kéo. Bác Việt đã gọi điện thống nhất với anh Nam là họ sẽ gặp nhau ở vị trí nào đó giữa bến Bính và thôn Hoành để hai người có mặt tại đó cùng lúc, không mất thời gian chờ nhau. Tìm vị trí hai người dự định gặp nhau, biết rằng vận tốc kéo xe của anh Nam là $5$ km/h và thuyền của bác Việt di chuyển với vận tốc $4$ km/h. Ngoài ra giả thiết rằng đường bờ biển từ thôn Hoành đến bến Bính là đường thẳng và bác Việt cũng luôn chèo thuyền tới một điểm trên bờ biển theo một đường thẳng.
	\loigiai{
		\immini{
			Ta mô hình hoá bài toán như trong hình vẽ. Trạm hải đăng ở vị tí $A$; bến Bính ở $B$ và thôn Hoành ở $C$. Giả sử bác Việt chèo thuyền cập bến ở vị trí $M$ và ta đặt $B M=x$ $(x>0)$. Để hai người không phải chờ nhau thì thời gian chèo thuyền bằng thời gian kéo xe nên ta có phương trình
		}{
			\begin{tikzpicture}[line cap=round, line join=round, font=\footnotesize, >=stealth, scale=0.5]
				\tikzset{label style/.style={font=\footnotesize}}
				\path (0,0) coordinate (B)
				(9.25,0) coordinate (C)
				(0,4) coordinate (A)
				(3,0) coordinate (M);
				\draw (M)--(A)--(B) node[midway, left]{$4$ km}--(C) (B)--(M) node[above, midway]{$x$ km};
				\draw[<->] ($(B)+(0,-0.5)$)--($(C)+(0,-0.5)$) node[below, midway]{$9,25$ km};
				\foreach\x/\y in {A/90,B/180,C/0,M/60} \draw[fill=black] (\x) circle (2pt) +(\y:0.6) node{$\x$};
				\draw pic[draw, angle radius=2mm]{right angle=A--B--M};
			\end{tikzpicture}	
			
		}
		\begin{align*}
			&\frac{\sqrt{x^{2}+16}}{4}=\frac{9{,}25-x}{5} \Rightarrow 25\left(x^2+16\right)=16\left(x-9{,}25\right)^2\\
			\Rightarrow\ &9x^2+296x-969=0 \Rightarrow \hoac{&x=3\\ &x=-\frac{323}{9}.}
		\end{align*}
		Thử lại, kết hợp với điều kiện $x>0$, ta được $BM=x=3$ km. Vậy vị trí gặp nhau cách thôn Bính $3$ km.
	}
\end{vd}

\begin{vd}
	\immini{
		Hằng ngày, bạn Hùng đều đón bạn Minh đi học tại một vị trí trên lề đường thẳng đến trường. Minh đứng tại vị trí $A$ cách lề đường một khoảng $50$ m để chờ Hùng. Khi nhìn thấy Hùng đạp xe đến địa điểm $B$, cách mình một đoạn
	}{
		\begin{tikzpicture}[line cap=round, line join=round, font=\footnotesize, >=stealth, scale=.6]
			\tikzset{label style/.style={font=\footnotesize}}
			\path (0,0) coordinate (H)
			(0,-2) coordinate (A)
			(-7.746,0) coordinate (B)
			(-1.849,0) coordinate (C);
			\draw (A)--(H) node[midway, right]{$50$ m}--(B)--(A) node[midway, sloped,below]{$200$ m}--(C) node[above,pos=0.3,sloped]{\tikz \draw[<-,thick] (0,0)--(1,0);};
			\foreach\x/\y in {A/-45,B/-90,C/90,H/90} \draw[fill=black] (\x) circle (1pt)+(\y:0.3) node{$\x$};
			\draw (B) node[above]{Hùng}
			($(B)+(1,0)$) node[above]{\tikz \draw[->,thick] (0,0)--(1,0);}
			(A) node[below left]{Minh};
			\draw pic[draw, angle radius=2mm]{right angle=A--H--B};
		\end{tikzpicture}	
		
	}
	$200$ m thì Minh bắt đầu đi bộ ra lề đường để bắt kịp xe. Vận tốc đi bộ của Minh là $5$ km/h, vận tốc xe đạp của Hùng là $15$ km/h. Hãy xác định vị trí $C$ trên lề đường để hai bạn gặp nhau mà không bạn nào phải chờ người kia (làm tròn kết quả đến hàng phần mười).
	\loigiai{
		Đặt $CH=x$ m, $x>0$. Ta có $BH=50\sqrt{15}$ m và $AC=\sqrt{x^2+2500}$ m. Tốc độ của Hùng và Minh lần lượt là $250$ m/phút và $\dfrac{250}{3}$ m/phút. Khi đó thời gian để Hùng và Minh di chuyển từ vị trí của mình đến $C$ lần lượt là $\dfrac{50\sqrt{15}-x}{250}$ phút và $\dfrac{3\sqrt{x^2+2500}}{250}$ phút. Để Hùng và Minh không phải chờ nhau thì
		\begin{align*}
			&\dfrac{50\sqrt{15}-x}{250}=\dfrac{3\sqrt{x^2+2500}}{250}\Rightarrow \left(x-50\sqrt{15}\right)^2=9\left(x^2+2500\right)\\
			\Rightarrow\ &8x^2+100\sqrt{15}x-15000=0 \Rightarrow \hoac{&x=-\frac{25\left(3\sqrt{7}+\sqrt{15}\right)}{4}\\ &x=\frac{25\left(3\sqrt{7}-\sqrt{15}\right)}{4}.}
		\end{align*}
		Thử lại, kết hợp với $x>0$, ta được $CH=\dfrac{25\left(3\sqrt{7}-\sqrt{15}\right)}{4}\approx 25{,}4$ m.
	}
\end{vd}

\baitaptl

\begin{bt}
	\immini{
		Có một nhà máy nước nọ muốn tìm vị trí $M$ để xây dựng trạm cấp nước sao cho khoảng cách từ nhà máy đến hai thị xã $B$, $C$ là bằng nhau. Biết hai thị xã trên lần lượt cách thành phố $A$ lần lượt $50$ km và $100$ km (như hình vẽ). Tính khoảng cách từ nhà máy nước đến thành phố $A$.
	}{
		\begin{tikzpicture}[line cap=round, line join=round, font=\footnotesize, >=stealth, scale=0.5]
			\tikzset{label style/.style={font=\footnotesize}}
			\path (0,0) coordinate (A)
			(10,0) coordinate (C)
			(0,5) coordinate (B)
			(4,0) coordinate (M);
			\draw (A)--(C) (M)--(B)--(A) node[midway, left]{$50$ km};
			\draw[<->] ($(A)+(0,-1)$) --($(C)+(0,-1)$) node[midway, below]{$100$ km};
			\foreach\x/\y in {A/-90,B/90,C/-90,M/60} \draw[fill=black] (\x) circle (1pt)+(\y:0.6) node{$\x$};
		\end{tikzpicture}	
	}
	\loigiai{
		Đặt $AM=x$ km, $x>0$. Khi đó, ta có
		\begin{align*}
			\sqrt{x^2+50^2}=BM=CM=100-x \Rightarrow x^2+2500=(100-x)^2\Rightarrow x=37{,}5.
		\end{align*}
		Thử lại, ta được khoảng cách từ nhà máy nước đến thành phố $A$ là $AM=37{,}5$ km. 
	}
\end{bt}

\begin{bt}
	\immini{
		Một người đứng ở điểm $A$ trên một bờ sông rộng $300$ m, chèo thuyền đến vị trí $D$, sau đó chạy bộ đến vị trí $B$ cách $C$ một khoảng $800$ m (hình vẽ). Vận tốc chèo thuyền là $6$ km/h, vận tốc chạy bộ là $10$ km/h và giả sử vận tốc dòng nước không đáng kể. Tính khoảng cách từ vị trí $C$ đến vị trí $D$, biết tổng thời gian người đó chèo thuyền và chạy bộ từ $A$ đến $B$ là $7{,}2$ phút.
	}{
		\begin{tikzpicture}[line cap=round, line join=round, font=\footnotesize, >=stealth, scale=0.4,rotate=90]
			\tikzset{label style/.style={font=\footnotesize}}
			\path (0,0) coordinate (A)
			(3,0) coordinate (C)
			(3,-8) coordinate (B)
			(3,-3.5) coordinate (D);
			\draw (D)--(A)--(C);
			\draw[shorten >=-0.5cm, shorten <=-0.5cm] (A)--($(A)-(C)+(B)$);
			\draw[shorten >=-0.5cm, shorten <=-0.5cm] (C)--(B);
			\draw[shorten >=-0.5cm, shorten <=-0.5cm] (A)--(C);
			\draw[shorten >=-0.5cm, shorten <=-0.5cm] (B)--($(A)-(C)+(B)$);
			\draw[<->] ($(A)+(0,0.4)$)--($(C)+(0,0.4)$) node[left, midway]{$300$ m};
			\draw[<->] ($(C)+(0.4,0)$)--($(B)+(0.4,0)$) node[above, midway]{$800$ m};
			\foreach\x/\y in {A/-135,B/-45,C/45,D/-150} \draw[fill=black] (\x) circle (2pt)+(\y:0.6) node{$\x$};
		\end{tikzpicture}	
		
	}
	\loigiai{
		Đặt $CD=x$ m, $x>0$. Ta có vận tốc của người đó khi trèo thuyền và chạy bộ lần lượt là $100$ m/phút và $\dfrac{500}{3}$ m/phút. Khi đó thời gian người đó chèo thuyền và chạy bộ lần lượt là $\dfrac{\sqrt{x^2+300^2}}{100}$ phút và $\dfrac{3\left(800-x\right)}{500}$ phút. Theo giả thiết, ta có phương trình
		\begin{align*}
			&\dfrac{\sqrt{x^2+300^2}}{100}+\dfrac{3\left(800-x\right)}{500}=7{,}2 \Rightarrow \sqrt{x^2+300^2}=240+\frac{3x}{5}\\
			\Rightarrow\ &x^2+300^2=\left(240+\frac{3x}{5}\right)^2\Rightarrow \frac{16x^2}{25}-288x+32400=0\Rightarrow x=225.
		\end{align*}
		Thử lại, ta được $CD=225$ m.
	}
\end{bt}

\begin{bt}
	\immini{
		Một con tàu biển $M$ rời cảng $O$ và chuyển động thẳng theo phương tạo với bờ biển một góc $60^{\circ}$. Trên bờ biển có hai đài quan sát $A$ và $B$ nằm về hai phía so với cảng $O$ và lần lượt cách cảng $O$ khoảng cách $1$ km và $2$ km.
	}{
		\begin{tikzpicture}[line cap=round, line join=round, font=\footnotesize, >=stealth, scale=0.6]
			\tikzset{label style/.style={font=\footnotesize}}
			\path (0,0) coordinate (O)
			(-2,0) coordinate (A)
			(4,0) coordinate (B)
			(60:3.5) coordinate (M);
			\draw (M)--(A)--(O) node[above,pos=0.6]{$1$ km}--(B) node[above, midway]{$2$ km}--(M)--(O) node[midway, right]{$x$ km};
			\foreach\x/\y in {A/-135,O/-90,M/90,B/-45} \draw[fill=black] (\x) circle (2pt)+(\y:0.5) node{$\x$};
			\draw pic[draw,angle radius=8mm]{angle=B--O--M}
			(25:0.6) node{$60^\circ$};
		\end{tikzpicture}	
		
	}
	\begin{enumerate}
		\item Đặt độ dài của $M O$ là $x \mathrm{~km}$. Biểu diễn khoảng cách từ tàu đến $A$ và từ tàu đến $B$ theo $x$.
		\item Tìm $x$ để khoảng cách từ tàu đến $B$ bằng $\dfrac{4}{5}$ khoảng cách từ tàu đến $A$.
		\item Tìm $x$ để khoảng cách từ tàu đến $B$ nhỏ hơn khoảng cách từ tàu đến $O$ đúng $500 \mathrm{~m}$.
	\end{enumerate}
	\loigiai{
		\begin{enumerate}
			\item Áp dụng định lí cô-sin, ta có
			\begin{align*}
				MA&=\sqrt{OM^2+OA^2-2OM\cdot OA\cdot \cos 120^\circ}=\sqrt{x^2+x+1},\\
				MB&=\sqrt{OM^2+OB^2-2OM\cdot OB\cdot \cos 60^\circ}=\sqrt{x^2-2x+4}.
			\end{align*}
			\item Theo yêu cầu, ta cần có $MB=\dfrac{4}{5}MA$, suy ra
			\begin{align*}
				&5\sqrt{x^2-2x+4}=4\sqrt{x^2+x+1} \Rightarrow 25\left(x^2-2x+4\right)=16\left(x^2+x+1\right)\\
				\Rightarrow\ & 9x^2-66x+84=0 \Rightarrow \hoac{&x=\frac{11-\sqrt{37}}{3}\\ &x=\frac{11+\sqrt{37}}{3}.}
			\end{align*}
			Thử lại, cả hai giá trị trên đều thỏa mãn. Vậy hoặc $x\approx 1{,}64$ km hoặc $x\approx 5{,}69$ km.
			\item Theo yêu cầu, ta cần có $MB=MO-0{,}5$, suy ra
			\begin{align*}
				&\sqrt{x^2-2x+4}=x-0{,}5 \Rightarrow x^2-2x+4=(x-0{,}5)^2\Rightarrow x=3{,}75.
			\end{align*}
			Thử lại, ta thấy thỏa mãn. Vậy $x=3{,}75$ km.
		\end{enumerate} 
	}
\end{bt}
% \subsubsection{Câu hỏi trắc nghiệm}
% \setcounter{ex}{0}
% \Opensolutionfile{ans}[ans/ans-PtQuyVeptBacHai]

% \begin{ex}%Bài 1. 
% 	\immini{Một con tàu biển $M$ rời cảng $O$ và chuyển động thẳng theo phương tạo với bờ biển một góc $60^{\circ}$. Trên bờ biển có hai đài quan sát $A$ và $B$ nằm về hai phía so với cảng $O$ và lần lượt cách cảng $O$ khoảng cách $1$ km và $2$ km (Hình bên).  Khoảng cách từ tàu đến $B$ bằng $\dfrac{4}{5}$ khoảng cách từ tàu đến $A$. Đặt độ dài của $MO$ là $x \mathrm{~km}$. Tìm $x$ biết ($x>2$) (Lưu ý:Làm tròn kết quả đến hàng phần trăm.)
% 		\choice
% 		{\True $5{,}69$ m}
% 		{$5{,}79$ m}
% 		{$5{,}89$ m}
% 		{$5{,}59$ m}}
% 	{\begin{tikzpicture}[scale=1, line cap=round, line join=round,font=\footnotesize, >=stealth]
% 			\def\a{1}
% 			\path 	(0:0) coordinate (O)
% 			++(0:2*\a) coordinate (B)
% 			++(180:3*\a) coordinate (A)
% 			($(O)!1.2!60:(B)$) coordinate (M);
% 			\draw[thick] (A)--(B)--(M)--cycle (O)--node[rotate=60, yshift=-1.5mm, scale=.7]{$x \mathrm{~km}$}(M);
% 			\foreach \i in {-1,-.9,...,2.1}{
% 				\draw (\i,0)--(\i-.1,-.1);
% 			}
% 			\foreach \x/\g in {A/200,B/-20,O/-90,M/90}
% 			\fill[black] (\x) circle (1pt)
% 			($(\g:3mm)+(\x)$) node[scale=.8] {$\x$};
% 			\draw (O) node[shift={(0.4,0.2)},scale=0.7]{$60^{\circ}$};
% 			%Hình chóp S.ABC có SA vuông góc đáy	
% 	\end{tikzpicture}}
% 	\loigiai{
% 		\begin{center}
% 			\begin{tikzpicture}[scale=1, line cap=round, line join=round,font=\footnotesize, >=stealth]
% 				\def\a{1}
% 				\path 	(0:0) coordinate (O)
% 				++(0:2*\a) coordinate (B)
% 				++(180:3*\a) coordinate (A)
% 				($(O)!1.2!60:(B)$) coordinate (M);
% 				\draw[thick] (A)--(B)--(M)--cycle (O)--node[rotate=60, yshift=-1.5mm, scale=.7]{$x \mathrm{~km}$}(M);
% 				\foreach \i in {-1,-.9,...,2.1}{
% 					\draw (\i,0)--(\i-.1,-.1);
% 				}
% 				\foreach \x/\g in {A/200,B/-20,O/-90,M/90}
% 				\fill[black] (\x) circle (1pt)
% 				($(\g:3mm)+(\x)$) node[scale=.8] {$\x$};
% 				%\draw pic[draw,double]{angle = B--O--M};
% 				\draw (O) node[shift={(0.4,0.2)},scale=0.7]{$60^{\circ}$};
% 				%Hình chóp S.ABC có SA vuông góc đáy	
% 			\end{tikzpicture}
% 		\end{center}
% 		\begin{enumerate}
% 			\item Áp dụng định lí hàm số Cosin trong tam giác, ta được:
% 			$$
% 			\begin{aligned}
% 				&AM^2=OM^2+OA^2-2OM\cdot OA\cos \left(180^\circ-60^\circ\right)
% 				=x^2+1+x \Rightarrow AM=\sqrt{x^2+x+1}. \\
% 				&BM^2=OM^2+OB^2-2OM\cdot OB\cos 60^\circ
% 				=x^2+4-2x \Rightarrow BM=\sqrt{x^2-2x+4}.
% 			\end{aligned}
% 			$$
% 			\item Theo đề ra ta có
% 			$$
% 			\begin{aligned}
% 				BM=\dfrac{4}{5}AM 
% 				&\Leftrightarrow \sqrt{x^2-2x+4}=\dfrac{4}{5}\sqrt{x^2+x+1} \\
% 				&\Leftrightarrow 25\left(x^2-2x+4\right)=16\left(x^2+x+1\right) \\
% 				&\Leftrightarrow 9x^2-66x+84=0 
% 				\Leftrightarrow \hoac{&x \approx 5{,}69\\
% 					&x\approx 1{,}64}{3}.
% 			\end{aligned}
% 			$$
			
% 		\end{enumerate}
% 	}
% \end{ex}
% %-------------------------------------------
% \begin{ex}%Bài 2. 
% 	\immini{Để leo lên một bức tường, bác Nam dùng một chiếc thang có chiều dài cao hơn bức tường đó $1$ m. Ban đầu, bác Nam đặt chiếc thang mà đầu trên của chiếc thang đó vừa chạm đúng vào mép trên bức tường (Hình 33a). Sau đó, bác Nam dịch chuyển chân thang vào gần chân tường thêm $0{,}5$ m thì bác Nam nhận thấy thang tạo với mặt đất một góc $60^{\circ}$ (Hình 33b). Bức tường cao bao nhiêu mét (làm tròn kết quả đến hàng phần mười)?
% 		\choice
% 		{\True $4{,}7$ m}
% 		{$5{,}7$ m}
% 		{$3{,}7$ m}
% 		{$4{,}4$ m}}
% 	{\begin{tikzpicture}[scale=0.8, font=\footnotesize,line join=round, line cap=round, >=stealth];
% 			\path
% 			(0,0) coordinate (O) node[below] {$C$}
% 			(2,5) coordinate (A)
% 			(0,5) coordinate (B)node[above] {$A$}
% 			(0.2,5.5) coordinate (B')
% 			(-2.5,0) coordinate (C)node[below] {$B$}
% 			(-1.8,0) coordinate (C')
% 			;
% 			\draw[very thick,color=orange,pattern=bricks] (O) rectangle (A);
% 			\draw[line width=4,blue] (B)--(C);
% 			\draw (O)--(C);
% 			\draw pic[draw,angle radius=2mm] {right angle = B--O--C'};
% 		\end{tikzpicture}
% 		\hfil
% 		\begin{tikzpicture}[scale=0.8, font=\footnotesize,line join=round, line cap=round, >=stealth];
% 			\path
% 			(0,0) coordinate (O) node[below] {$G$}
% 			(2,5) coordinate (A)
% 			(0,5) coordinate (B)node[left] {$D$}
% 			(0.2,5.5) coordinate (B')
% 			(-2.5,0) coordinate (C)
% 			(-1.8,0) coordinate (C')node[below] {$E$}
% 			;
% 			\draw[very thick,color=orange,pattern=bricks] (O) rectangle (A);
% 			\draw[line width=4,blue] (B')--(C');
% 			\draw (O)--(C);
% 			%\node[rectangle,name=node1,align=none,left] at ($ (C)!0.5!(B) $){$4$m};
% 			\pic["$60^\circ$"{shift={(1pt,0pt)}},draw,angle radius=4mm,angle eccentricity=1.79] {angle =O--C'--B'};
% 			\draw pic[draw,angle radius=2mm] {right angle = B'--O--C'};
% 		\end{tikzpicture}
% 	}
% 	\loigiai{
% 		Gọi $x$ là chiều cao của bức tường. Điều kiện $x>0$, đơn vị: mét. \\
% 		Suy ra chiều dài của thang là $x+1$ m. \\
% 		Áp dụng định lí Py-ta-go trong hình 33a, ta được
% 		$$BC^2=(x+1)^2-x^2=2x+1 \Rightarrow BC=\sqrt{2x+1} \mathrm{~m}.$$
% 		Áp dụng hệ thức lượng trong tam giác vuông ở hình 33b, ta được
% 		$$
% 		\begin{aligned}
% 			DG=EG\tan 60 &\Leftrightarrow x=\sqrt{3}\left(\sqrt{2x+1}-0{,}5\right) \\
% 			&\Leftrightarrow \sqrt{2x+1}=\dfrac{x}{\sqrt{3}}+0{,}5 
% 			\Leftrightarrow 2x+1=\left(\dfrac{x}{\sqrt{3}}+0{,}5\right)^2 \\
% 			&\Leftrightarrow \dfrac{x^2}{3}+\left(\dfrac{1}{\sqrt{3}}-2\right)x-\dfrac{3}{4}=0
% 			\Leftrightarrow \hoac{&x\approx 4{,}74 \text{ (nhận)} \\ &x\approx -0{,}74 \text{ (loại)}.}
% 		\end{aligned}$$
% 	}
% \end{ex}
% %-------------------------------------------
% \begin{ex}%Bài 3.
% 	\immini{
% 		Một người đứng ở điểm $A$ trên một bờ sông rộng $300$ m, chèo thuyền đến vị trí $D$, sau đó chạy bộ đến vị trí $B$ cách $C$ một khoảng $800$ m như Hình 34.  Vận tốc chèo thuyền là $6$ km/h, vận tốc chạy bộ là $10$ km/h và giả sử vận tốc dòng nước không đáng kể. Tính khoảng cách từ vị trí $C$ đến $D$, biết tổng thời gian người đó chèo thuyền và chạy bộ từ $A$ đến $B$ là $7{,}2$ phút.
% 		\choice
% 		{$224$ m}
% 		{\True $225$ m} 
% 		{$226$ m} 
% 		{$227$ m}}
% 	{\begin{tikzpicture}[scale=1, line cap=round, line join=round, font=\footnotesize, >=stealth]
% 			\path (0,0) coordinate (A)
% 			++(0:2) coordinate (C)
% 			++(-90:2) coordinate (D)++(-90:2)coordinate (B)
% 			;
% 			\draw (A)--node[yshift=1.5mm,scale=.7]{$300$m}(C)--(B)--cycle
% 			(0,-4)--(A)--(D);
% 			\draw[<->,dashed] ([xshift=3mm]C)--node[rotate=90,scale=.7, yshift=-2mm]{$800$m}([xshift=3mm]B);
% 			\foreach \i in {0,-.1,...,-4.1}{
% 				\draw (0,\i)--(-.1,\i-0.1);
% 			}
% 			\foreach \i in {0,-.1,...,-4.1}{
% 				\draw (2,\i)--(2.1,\i-0.1);
% 			}
% 			\foreach \i/\j in {A/160,C/20,D/200,B/200}{
% 				\fill[black] (\i)circle(.7pt)
% 				($(\i)+(\j:2mm)$)node[scale=.8]{$\i$};
% 			}
% 		\end{tikzpicture}
% 	}
% 	\loigiai{
% 		Đặt $CD=x\in (0;0{,}8)$, đơn vị kilomet. Suy ra $BD=0{,}8-x$ km. \\
% 		Khoảng cách từ $A$ đến $D$ là $AD=\sqrt{0{,}3^2+x^2}$ km. \\
% 		Thời gian đi từ $A$ đến $D$ là $t_{AD}=\dfrac{\sqrt{0{,}3^2+x^2}}{6}$ h.\\
% 		Thời gian đi từ $D$ đến $B$ là $t_{DB}=\dfrac{0{,}8-x}{10}$ h.\\
% 		Theo đề ra, ta có phương trình $\dfrac{\sqrt{0{,}3^2+x^2}}{6}+\dfrac{0{,}8-x}{10}=\dfrac{7{,}2}{60}$. \hfill $(1)$ \\
% 		Giải $(1)$ ta được $x=\dfrac{9}{40}$ km = $225$ m. \\
% 		Vậy khoảng cách từ $C$ đến $D$ là $225$ m.
% 	}
% \end{ex}
% %-------------------------------------------
% \begin{ex}%Bài 4.
% 	Một ngọn hải đăng đặt tại vị trí $A$ cách bờ biển một khoảng cách $AB=4$km. Trên bờ biển có một cái kho ở vị trí $C$ cách $B$ một khoảng là $7$ km. Người canh hải đăng có thể chèo thuyền từ $A$ đến vị trí $M$ trên bờ biển với vận tốc $3$ km/h rồi đi bộ đến $C$ với vận tốc $5$ km/h như hình sau.  Tính khoảng cách từ vị trí $B$ đến $M$, biết thời gian người đó đi từ $A$ đến $C$ là $148$ phút.
% 	\begin{center}
% 		\begin{tikzpicture}[>=stealth,line join=round,line cap=round,font=\footnotesize,scale=0.8]
% 			\path (0,0) coordinate (B) (0,4) coordinate (A) (7,0) coordinate (C)
% 			($(B)!2/5!(C)$) coordinate (M);
% 			\draw (A)--(B)--(C)--cycle (A)--(M);
% 			\draw (-1,0)--(B) (C)--(8,0);
% 			\foreach \x/\g in {A/90,B/135,C/60,M/45}
% 			\draw[fill=black] (\x) circle (1pt) + (\g:.3)node{$\x$};
% 			\draw[dashed] (B)--++(0,-0.5) (C)--++(0,-0.5);
% 			\draw[<->] (0,-0.3)--(7,-0.3);
% 			\draw (3.5,-0.3) node[below]{$7$km};
% 			\draw (0,2) node[left]{$4$km};
% 		\end{tikzpicture}
% 		\hspace*{2cm}
% 		\begin{tikzpicture}[smooth,scale=.5,line join=round,scale=0.4]
% 			\draw [fill=black,smooth,line join=round,black] (-6.14475,0)..controls +(60:.5) and +(120:.5)..(-5.727,0.3825)..controls +(60:.5) and +(120:.5)..(-4.1775,0.9825)..controls +(60:.5) and +(120:.5)..(-3.0275,2)..controls +(60:.5) and +(120:.5)..(-2.,2.)--(-1.8575,4.5325)--(-1.6275,3.1825)..controls +(70:1) and +(-110:1)..(1.8,5.5825)--(2.,2.)..controls +(60:.5) and +(120:.5)..(2.9225,1.9325)..controls +(60:.5) and +(120:.5)..(4.2225,1.0825)..controls +(60:.5) and +(120:.5)..(5.9225,0.4325)..controls +(60:.5) and +(120:.5)..(6.7725,0)--cycle;
% 			\draw [fill=black,smooth,black] (-1.77,5.8825)--(-1.62,8.7325)..controls +(70:1) and +(-110:1)..(1.48,10.8325)--(1.612,8.0825)..controls +(-110:1) and +(70:1)..(-1.77,5.8825);
% 			\draw [fill=black,black] (-1.45,11.4325)--(-1.22,15) arc (180:0:{1.22} and {.5})--(1.22,15)--(1.34,13.3325)..controls +(-110:1) and +(70:1)..(-1.45,11.4325);
% 			\draw [line width=.1cm,black] (-1.18,15)--(-1.22,15.6) arc (180:0:{1.22} and {.5})--(1.22,15.6)--(1.18,15);
% 			\draw[fill=black,black] (0,14) ellipse ({2} and {.5});
% 			\draw[fill=black,black] (0,17.5) ellipse ({.2} and {.8});
% 			\draw [fill=black,black] (0,16.4) circle (.35);
% 			\draw[fill=black,black] (-.9,17) arc (180:0:{.9 and .9});
% 			\draw [line width=.08cm,black] (0.87,15)--(0.87,16);
% 			\draw [line width=.08cm,black] (-0.87,15)--(-0.87,16);
% 			\draw [line width=.08cm,black] (0.33,15)--(0.33,16.1);
% 			\draw [line width=.08cm,black] (-0.33,15)--(-0.33,16.1);
% 			\draw [line width=.06cm,black] (0.81,16)--(0.81,17.2);
% 			\draw [line width=.06cm,black] (-0.81,16)--(-0.81,17.2);
% 			\path 
% 			(-1,5.3) coordinate (A)
% 			(.8,11.5) coordinate (B);
% 			\draw[fill=black,black] ($(A)+(60:.6)$)--($(A)+(120:.6)$)--($(A)+(-120:.6)$)--($(A)+(-60:.6)$)--cycle;
% 			\draw[fill=black,black] ($(B)+(60:.6)$)--($(B)+(120:.6)$)--($(B)+(-120:.6)$)--($(B)+(-60:.6)$)--cycle;
% 			%\draw[yellow](-2,2)--(-1.22,15.9)--(1.22,15.9)--(2,2)--cycle; 
% 		\end{tikzpicture}
% 	\end{center}
% 	\choice
% 	{$1$ km}
% 	{$2$ km} 
% 	{\True $3$ km} 
% 	{$4$ km} 
% 	\loigiai{
% 		Đặt $BM=x \in (0;7)$, đơn vị km. Suy ra $MC=7-x$ km. \\
% 		Ta có $AM=\sqrt{x^2+4^2}=\sqrt{x^2+16}$ km. \\
% 		Thời gian đi từ $A$ đến $M$	là $t_{AM}=\dfrac{\sqrt{x^2+16}}{3}$ h. \\
% 		Thời gian đi từ $M$ đến $C$	là $t_{MC}=\dfrac{7-x}{5}$ h. \\
% 		Theo đề ra ta có phương trình $\dfrac{\sqrt{x^2+16}}{3}+\dfrac{7-x}{5}=\dfrac{148}{60}$. \hfill $(1)$ \\
% 		Giải $(1)$ ta được $x=3$ km. \\
% 		Vậy khoảng cách từ $B$ đến $M$ là $3$ km.
% 	}
% \end{ex}
% %-------------------------------------------
% \begin{ex}%Bài 5.
% 	\immini{Bác Việt sống và làm việc tại trạm hải đăng (vị trí A) cách bờ biển $4$ km. Hàng tuần bác chèo thuyền vào vị trí gần nhất trên bờ biển là bến Bính (vị trí B) để nhận hàng hoá do cơ quan cung cấp. Tuần này, do trục trặc về vận chuyển nên toàn bộ số hàng vẫn đang nằm ở thôn Hoành (vị trí C), bên bờ biển cách bến Bính $9{,}25$ km và sẽ được anh Nam vận chuyển trên con đường dọc bờ biển tới bến Bính bằng xe kéo. Bác Việt đã gọi điện thống nhất với anh Nam là họ sẽ gặp nhau ở vị trị nào đó giữa bến Bính và thôn Hoành để hai ngươi có mặt tại đó cùng lúc, không mất thời gian chờ nhau.}
% 	{\begin{tikzpicture}[scale=1, line cap=round, line join=round, font=\footnotesize, >=stealth]
% 			\path (0,0)coordinate (A)
% 			++(-90:2)coordinate (B)
% 			++(0:1.5)coordinate (M)++(0:3)coordinate (C)
% 			;
% 			\draw (A)--node[rotate=90,yshift=1.5mm,scale=.7]{$4$km}(B)--node[yshift=1.5mm,scale=.7]{$x$}(M)--(C)
% 			(A)--(M);
% 			\draw[<->,dashed] ([yshift=-3mm]B)--node[yshift=-1.5mm,scale=.7]{$9{,}25$km}([yshift=-3mm]C);
% 			\foreach \a/\b in {A/90,B/180,M/40,C/30}
% 			\fill[black] (\a)circle(.7pt) ($(\a)+(\b:2mm)$)node[scale=.8]{$\a$};
% 	\end{tikzpicture}}
% 	Tìm vị trí hai người dự định gặp nhau (vị trí M), biết rằng vận tốc kéo xe của anh Nam là $5$ km/h và thuyền của bác Việt di chuyển với vận tốc $4$ km/h. Ngoài ra giả thiết rằng đường bờ biển từ thôn Hoành đến bến Bính là đường thẳng và bác Việt cũng luôn chèo thuyền tới một điểm trên bờ biển theo một đường thẳng. Khi đó khoảng cách $BM$ bằng
% 	\choice
% 	{$1$ km}
% 	{$2$ km} 
% 	{\True $3$ km} 
% 	{$4$ km} 
% 	\loigiai{
% 		\immini{
% 			Ta mô hình hoá bài toàn như trong hình bên: Trạm hải đăng ở vị trí $A$; bến Bính ở $B$ và thôn Hoành ở $C$. Giả sử bác Việt chèo thuyền cập bến ở vị trí $M$ và ta đặt $BM=x\,\, (0<x<9{,}25)$. Để hai người không phải chờ nhau thì thời gian chèo thuyển bằng thời gian kéo xe nên ta có phương trình:
% 		}
% 		{
% 			\begin{tikzpicture}[scale=1, line cap=round, line join=round, font=\footnotesize, >=stealth]
% 				\path (0,0)coordinate (A)
% 				++(-90:2)coordinate (B)
% 				++(0:1.5)coordinate (M)++(0:3)coordinate (C)
% 				;
% 				\draw (A)--node[rotate=90,yshift=1.5mm,scale=.7]{$4$km}(B)--node[yshift=1.5mm,scale=.7]{$x$}(M)--(C)
% 				(A)--(M);
% 				\draw[<->,dashed] ([yshift=-3mm]B)--node[yshift=-1.5mm,scale=.7]{$9{,}25$km}([yshift=-3mm]C);
% 				\foreach \a/\b in {A/90,B/180,M/40,C/30}
% 				\fill[black] (\a)circle(.7pt) ($(\a)+(\b:2mm)$)node[scale=.8]{$\a$};
% 			\end{tikzpicture}
% 		}
% 		$$\dfrac{\sqrt{x^{2}+16}}{4}=\dfrac{9{,}25-x}{5}.$$
% 		Giải phương trình này sẽ tìm được $x=3$ (nhận); $x=-\dfrac{323}{9}$ (loại).\\
% 		Vậy hai người gặp nhau ở vị trí $M$ cách bến Bính $B$ là $3$ kilomet.
% 	}
% \end{ex}
% %-------------------------------------------
% \begin{ex}%Bài 6.
% 	Cho tứ giác $ABCD$ có $AB\perp CD;\; AB=2;\; BC=13;\; CD=8; DA=5$ (H.6.21). Gọi $H$ là giao điểm của $AB$ và $CD$ và đặt $x=AH$. Hãy thiết lập một phương trình để tính độ dài $x$, từ đó tính diện tích tứ giác $ABCD$ bằng
% 	\begin{center}
% 		\begin{tikzpicture}[scale=1, line cap=round, line join=round, font=\footnotesize, >=stealth]
% 			\def\a{2} \pgfmathsetmacro\b{2+\a}
% 			\pgfmathsetmacro\c{sqrt(13^2-(\b)^2)}
% 			\pgfmathsetmacro\d{sqrt(5^2-(\a)^2)}
% 			\path (0,0) coordinate (H)
% 			(0,-\a)coordinate (A)
% 			(0,-\b)coordinate (B)
% 			(\c,0) coordinate (C)
% 			(\d,0) coordinate (D)
% 			;
% 			\draw pic[draw, angle radius=2mm,fill=gray!20]{right angle=A--H--D};
% 			\draw[dashed] (A)--node[rotate=90,yshift=1.5mm,scale=.8]{$x$}(H)--(D);
% 			\draw (A)--node[rotate=90,yshift=1.5mm,scale=.8]{$2$}(B)--node[rotate=20,yshift=1.5mm,scale=.8]{$13$}(C)--node[yshift=1.5mm,scale=.8]{$8$}(D)--node[rotate=35,yshift=1.5mm,scale=.8]{$5$}cycle;
% 			\foreach \a/\b in {H/120,A/180,B/-90,C/0,D/90}
% 			\fill[black] (\a) circle(.7pt)
% 			($(\a)+(\b:2mm)$)node[scale=.8]{$\a$};
% 		\end{tikzpicture}
% 	\end{center}
% 	\choice
% 	{$23$ m}
% 	{\True $24$ m}
% 	{$25$ m}
% 	{$26$ m}
% 	\loigiai{
		
% 		Áp dụng định lí Py-ta-go trong tam giác vuông $AHD$, ta được
% 		$$HD^2=AD^2-AH^2=25-x^2 \Rightarrow HD=\sqrt{25-x^2}.$$
% 		Áp dụng định lí Py-ta-go trong tam giác vuông $BHC$, ta được
% 		$$
% 		\begin{aligned}
% 			HB^2+HC^2=BC^2 &\Leftrightarrow (x+2)^2+\left(8+\sqrt{25-x^2}\right)^2=169 \\
% 			&\Leftrightarrow 4\sqrt{25-x^2}+x-19=0 \qquad \hfill  (1) \\.
% 		\end{aligned}$$
% 		Giải $(1)$ ta được $x=3$ (nhận); $x=-\dfrac{13}{17}$ (loại). \\
% 		Vậy $AH=3$; $HD=4$. Suy ra $S_{ABCD}=S_{BHC}-S_{AHD}= 24$. 
% 	}
% \end{ex}
% %-------------------------------------------
% \begin{ex}%Bài 7.
% 	Hằng ngày bạn Hùng đều đón bạn Minh đi học tại một vị trí trên lề đường thẳng đến trường. Minh đứng tại vị trí $A$ cách lề đường một khoảng $50$ m để chờ Hùng. Khi nhìn thấy Hùng đạp xe đến địa điểm $B$, cách mình một đoạn $200$ m thì Minh bắt đầu đi bộ ra lề đường để bắt kịp xe. Vận tốc đi bộ của Minh là $5 \mathrm{~km}/\mathrm{h}$, vận tốc xe đạp của Hùng là $15 \mathrm{~km}/\mathrm{h}$. Hãy xác định vị trí $C$ trên lề đường (hình sau) để hai bạn gặp nhau mà không bạn nào phải chờ người kia (làm tròn kết quả đến hàng phần mười). Khi đó đoạn thẳng $CH$ bằng
% 	\begin{center}
% 		\begin{tikzpicture}[scale=0.9, font=\footnotesize,line join=round, line cap=round, >=stealth];
% 			\path
% 			(0,0) coordinate (B) node[above] {Hùng}
% 			(6,-2) coordinate (A)node[below] {Minh}
% 			(6,0) coordinate (H)
% 			(4,0) coordinate (C)								
% 			;													
% 			\path
% 			($(A)!.5!(B)$)coordinate (E)node[left] {$200$m}
% 			($(A)!.6!(H)$)coordinate (M)node[right] {$50$m};		
% 			\draw (A)--(B)--(H)--(A)--(C);							
% 			\foreach \x/\g in{A/0,B/-90,C/90,H/90}
% 			\fill[black](\x)circle(1pt) ($(\x)+(\g:3mm)$)node{$\x$};
% 			\draw pic[draw,angle radius=2mm] {right angle = A--H--B};		
% 		\end{tikzpicture}
% 	\end{center}
% 	\choice
% 	{\True $25{,}4$ m}
% 	{$25{,}6$ m}
% 	{$25{,}7$ m}
% 	{$25{,}2$ m}
% 	\loigiai{
% 		Đổi đơn vị $200 \mathrm{m}=0{,}2\mathrm{km}$; $50 \mathrm{m}=0{,}05\mathrm{km}$. \\
% 		Áp dụng định lí Py-ta-go trong tam giác vuông $ABH$ ta được
% 		$$BH^2=AB^2-AH^2=\dfrac{3}{80} \Rightarrow BH=\dfrac{\sqrt{15}}{20}.$$
% 		Đặt $HC=x \in \left(0; \dfrac{\sqrt{15}}{20}\right)$, đơn vị $\mathrm{km}$. \\
% 		Ta có $AC=\sqrt{x^2+0{,}05^2} \Rightarrow t_{AC}=\dfrac{\sqrt{x^2+0{,}05^2}}{5}$; $t_{BC}=\dfrac{\dfrac{\sqrt{15}}{20}-x}{15}$. \\
% 		Để Minh và Hùng không phải chờ nhau thì ta phải có
% 		$$
% 		\begin{aligned}
% 			t_{AC}=t_{BC} &\Leftrightarrow \dfrac{\sqrt{x^2+0{,}05^2}}{5} = \dfrac{\dfrac{\sqrt{15}}{20}-x}{15} \\
% 			&\Leftrightarrow 8x^2=\dfrac{\sqrt{15}}{20}x-\dfrac{3}{200} =0
% 			\Leftrightarrow \hoac{&x\approx 0{,}0254 \text{ (nhận)} \\ &x\approx -0{,}0738 \text{ (loại)}.}
% 		\end{aligned}
% 		$$
% 		Vậy $C$ cách $H$ một khoảng gần bằng $25{,}4$ mét.
% 	}
% \end{ex}
% %-------------------------------------------

% \Closesolutionfile{ans}
% \begin{indapan}{10}
% 	{ans/ans-PtQuyVeptBacHai}
% \end{indapan}